\documentclass[10pt,UTF8]{article}

\usepackage{geometry}

\usepackage{ctex}

\geometry{a4paper,scale=0.9}
\ctexset{today=old}

\usepackage[bookmarksnumbered=true]{hyperref}  % 生成大纲



\usepackage{times}  % 设置正文字体为 Adobe Times Roman

% \usepackage{makecell}
% \usepackage{multirow} % 表格多行合并

% \usepackage{graphicx} %插入图片的宏包
% \usepackage{subfigure}
% \usepackage{float} %设置图片浮动位置的宏包

\usepackage{amsmath}
\usepackage{amssymb}
\usepackage{mathtools}

\usepackage{physics}

% \usepackage{siunitx}
% % 关闭警告
% \ExplSyntaxOn
% \msg_redirect_name:nnn { siunitx } { physics-pkg } { none }
% \ExplSyntaxOff

\allowdisplaybreaks[3]  % 允许换页


%%%%%%%%%%%%%%%%%%%%%%%% tcolorbox %%%%%%%%%%%%%%%%%%%%%%%%%%
\usepackage[most]{tcolorbox}

\newenvironment{definition box}
{\begin{tcolorbox}[colback=blue!5!white,colframe=blue!75!black,parbox=false,before upper=\par,before lower=\par]}
{\end{tcolorbox}}

\newenvironment{theorem box}
{\begin{tcolorbox}[colback=yellow!5!white,colframe=yellow!75!black,parbox=false,before upper=\par,before lower=\par]}
{\end{tcolorbox}}
%%%%%%%%%%%%%%%%%%%%%%%%%%%%%%%%%%%%%%%%%%%%%%%%%%%%%%%%%%%%
            
%%%%%%%%%%%%%%%%%%%%% amsthm %%%%%%%%%%%%%%%%
\usepackage{amsthm}

\theoremstyle{definition}
\newtheorem{definition inner}{Definition}[section]
\newenvironment{definition}[1][]
{\begin{definition box}\begin{definition inner}[#1]\hfill\par}
{\end{definition inner}\end{definition box}}

\theoremstyle{plain}
\newtheorem{theorem inner}{Theorem}[section]
\newenvironment{theorem}[1][]
{\begin{theorem box}\begin{theorem inner}[#1]\hfill\par}
{\end{theorem inner}\end{theorem box}}

\theoremstyle{definition}
\newtheorem*{proof inner}{\textit{Proof}}
\renewenvironment{proof}[1][]
{\begin{proof inner}[#1]\hfill\par}
{\qed\end{proof inner}}

\theoremstyle{remark}
\newtheorem*{remark}{\textit{Remark}}
%%%%%%%%%%%%%%%%%%%%%%%%%%%%%%%%%%%%%%%%%%%%%%


\newcommand{\mycases}[2][0.8]{
    \left\{\vphantom{\scalebox{#1}{
        $\begin{matrix}#2\end{matrix}$
    }}\right.\!\!\!\!
    \begin{array}{l@{\quad}l}#2\end{array}
}

\newcommand{\ju}{\mathrm{j}}  % 虚数单位
\newcommand{\e}{\mathrm{e}{\mkern 1 mu}}

\newcommand{\Ev}{\,\mathcal{E}\ensuremath{v}\qty}
\newcommand{\Od}{\,\mathcal{O}\ensuremath{d}\qty}

\renewcommand{\Re}{\,\mathcal{R}\ensuremath{e}\qty}
\renewcommand{\Im}{\,\mathcal{I}\ensuremath{m}\qty}

\newcommand{\sinc}{\operatorname{sinc}} 
\newcommand{\rect}{{\mkern 1.5 mu}\operatorname{rect}\!\qty} % 矩形函数


\begin{document}

\part*{The z-Transform}

\section{The z-Transform}

\begin{definition}[The z-Transform]
The (bilateral) z-transform of a general signal $x[n]$ is defined as
\begin{equation*}
    X(z) =
    \sum_{n=-\infty}^{+\infty} x[n] z^{-n}
\end{equation*}
where $z=r\e^{\ju\omega}$ is a complex variable.
\end{definition}

\begin{remark}
\textbf{Connection with the Fourier Transform}
\begin{equation*}
    \mathcal{Z}\{x[n]\} = X(z) =
    X(r\e^{\ju\omega}) = \mathcal{F}\{x[n] r^{-n}\}
\end{equation*}
\end{remark}

\begin{definition}[The Inverse z-Transform]
The inverse z-transform of a function $X(z)$ is defined as
\begin{equation*}
    x[n] = \frac{1}{2\pi\ju} \oint_{\abs{z}=r}\!\!\!\!\!\!\!\!
    X(z) z^{n-1} \dd{z}
\end{equation*}
A contour integration around a counterclockwise closed
circular contour centered at the origin and with radius $r$.
\end{definition}



\section{Properties of the ROC}

\begin{enumerate}
    \item The ROC consists of circular regions centered about the origin.
    \item The ROC dose not include any poles.
    \item If $x[n]$ is of finite duration and is absolutely integrable,
    then the ROC is the entire $z$-plane.
    \item If $x[n]$ is right-sided, then the ROC is of the form $\abs{z} > r_0$
    for some $r_0$.
    \item If $x[n]$ is left-sided, then the ROC is of the form $0<\abs{z}<r_0$
    for some $r_0$.
    \item If $x[n]$ is two-sided, then the ROC is of
    the form $r_1 < \Re(z) < r_2$ for some $r_1$ and $r_2$.
\end{enumerate}



\section{Basic z-Transform Pairs}

\begin{center}
\begin{tabular}
    { c @{\qquad} c @{\qquad} c}
    Signal & Laplace Transform & ROC \\
    \hline $\displaystyle
    \delta[n]
    $ & $\displaystyle
    1
    $ & $
    \mathbb{C} \vphantom{\mycases{1\\1}}
    $ \\\hline $
    \delta(n-m)
    $ & $\displaystyle
    z^{-m}
    $ & $
    \mathbb{C} \vphantom{\mycases{1\\1}}
    $ \\\hline $
    u[n]
    $ & $\displaystyle
    \frac{1}{1-z^{-1}}
    $ & $
    \abs{z} > 1 \vphantom{\mycases{1\\1}}
    $ \\\hline $
    -u[-n-1]
    $ & $\displaystyle
    \frac{1}{1-z^{-1}}
    $ & $
    \abs{z} < 1 \vphantom{\mycases{1\\1}}
    $ \\\hline $
    a^n u[n]
    $ & $\displaystyle
    \frac{1}{1-a z^{-1}}
    $ & $
    \abs{z} > \abs{a} \vphantom{\mycases{1\\1}}
    $ \\\hline $
    -a^n u[-n-1]
    $ & $\displaystyle
    \frac{1}{1-a z^{-1}}
    $ & $
    \abs{z} < \abs{a} \vphantom{\mycases{1\\1}}
    $ \\\hline $
    b^{\abs{n}}
    $ & $\displaystyle
    \frac{1}{1-b z^{-1}} - \frac{1}{1-b^{-1} z^{-1}}
    $ & $
    \abs{b} < \abs{z} < \abs{b}^{-1} \vphantom{\mycases{1\\1}}
    $ \\\hline $\displaystyle
    n a^n u[n]
    $ & $\displaystyle
    \frac{a z^{-1}}{(1-a z^{-1})^2}
    $ & $
    \abs{z} > \abs{a} \vphantom{\mycases{1\\1}}
    $ \\\hline $
    -n a^n u[-n-1]
    $ & $\displaystyle
    \frac{a z^{-1}}{(1-a z^{-1})^2}
    $ & $
    \abs{z} < \abs{a} \vphantom{\mycases{1\\1}}
    $ \\\hline $\displaystyle
    \cos(\omega_0 n) u[n]
    $ & $\displaystyle
    \frac{1 - z^{-1}\cos\omega_0}{1 - 2 z^{-1}\cos\omega_0 + z^{-2}}
    $ & $\displaystyle
    \abs{z} > 1 \vphantom{\mycases{1\\1}}
    $ \\\hline $\displaystyle
    \sin(\omega_0 n) u[n]
    $ & $\displaystyle
    \frac{z^{-1}\sin\omega_0}{1 - 2 z^{-1}\cos\omega_0 + z^{-2}}
    $ & $\displaystyle
    \abs{z} > 1 \vphantom{\mycases{1\\1}}
    $ \\\hline $\displaystyle
    r^n\cos(\omega_0 n) u[n]
    $ & $\displaystyle
    \frac{1 - z^{-1}r\cos\omega_0}{1 - 2 z^{-1}r\cos\omega_0 + z^{-2}r^2}
    $ & $\displaystyle
    \abs{z} > r \vphantom{\mycases{1\\1}}
    $ \\\hline $\displaystyle
    r^n\sin(\omega_0 n) u[n]
    $ & $\displaystyle
    \frac{z^{-1}r\sin\omega_0}{1 - 2 z^{-1}r\cos\omega_0 + z^{-2}r^2}
    $ & $\displaystyle
    \abs{z} > r \vphantom{\mycases{1\\1}}
    $ \\\hline
\end{tabular}
\end{center}



\section{Properties of the Bilateral z-Transform}

\subsection{}

\begin{center}
\begin{tabular}
    { c @{\qquad} c @{\qquad} c}
    Signal & Transform & ROC \\
    \hline $\displaystyle
    x[n]
    $ & $\displaystyle
    X(z)
    $ & $
    R
    $ \\\hline $
    x_1[n]
    $ & $\displaystyle
    X_1(z)
    $ & $
    R_1
    $ \\\hline $
    x_2[n]
    $ & $\displaystyle
    X_2(z)
    $ & $
    R_2
    $ \\\hline $
    a x_1[n] + b x_2[n]
    $ & $\displaystyle
    a X_1(z) + b X_2(z)
    $ & $
    \text{At least } R_1 \cap R_2
    $ \\\hline $
    x[n-n_0]
    $ & $\displaystyle
    z^{-n_0} X(z)
    $ & $
    R
    $ \\\hline $\displaystyle
    x[-n]
    $ & $\displaystyle
    X\qty(\frac{1}{z})
    $ & $\displaystyle
    \frac{1}{R} \vphantom{\mycases{1\\1}}
    $ \\\hline $\displaystyle
    z_0^n x[n]
    $ & $\displaystyle
    X\qty(\frac{z}{z_0})
    $ & $
    \abs{z_0} R \vphantom{\mycases{1\\1}}
    $ \\\hline $\displaystyle
    x_{(k)}[n] = \mycases{
        x[n/k] \\ 0
    }
    $ & $\displaystyle
    X(z^k)
    $ & $
    R^{1/k}
    $ \\\hline $
    x^*[n]
    $ & $\displaystyle
    X^*(z^*)
    $ & $
    R
    $ \\\hline $
    x_1[n] \ast x_2[n]
    $ & $\displaystyle
    X_1(z) X_2(z)
    $ & $
    \text{At least } R_1 \cap R_2
    $ \\\hline $\displaystyle
    n x[n]
    $ & $\displaystyle
    -z \dv{z} X(z)
    $ & $\displaystyle
    R \vphantom{\mycases{1\\1}}
    $ \\\hline $\displaystyle
    \sum_{k=-\infty}^{n} x[k]
    $ & $\displaystyle
    \frac{1}{1-z^{-1}} X(z)
    $ & $\displaystyle
    \text{At least } R \cap \{\abs{z}>1\} \vphantom{\mycases{1\\1}}
    $ \\\hline
\end{tabular}
\end{center}

\subsection{Initial- and Final-Value Theorem}

\begin{theorem}[Initial-Value and Final-Value Theorem]
    If $x[n]=0$ for $n<0$, then
    \begin{gather*}
        x(0) = \lim_{z\to+\infty} X(z) \\
        \lim_{n\to+\infty} x[n] = \lim_{z\to 1} (z-1) X(z)
    \end{gather*}
\end{theorem}



\section{The Unilateral z-Transform}

\subsection{}

\begin{definition}[The Unilateral z-Transform]
The unilateral z-transform of a general signal $x[n]$ is defined as
\begin{equation*}
    \mathfrak{X}(z) =
    \sum_{n=0}^{+\infty} x[n] z^{-n}
\end{equation*}
where $z=r\e^{\ju\omega}$ is a complex variable.
\end{definition}


\subsection{Properties of the Unilateral z-Transform}

For causal signals, the unilateral z-transform is the same as the bilateral z-transform.

\subsubsection{}

\begin{center}
\begin{tabular}
    { c @{\qquad} c @{\qquad} c}
    Signal & Transform \\
    \hline $\displaystyle
    x[n]
    $ & $\displaystyle
    \mathfrak{X}(z)
    $ \\\hline $
    x_1[n]
    $ & $\displaystyle
    \mathfrak{X}_1(z)
    $ \\\hline $
    x_2[n]
    $ & $\displaystyle
    \mathfrak{X}_2(z)
    $ \\\hline $
    a x_1[n] + b x_2[n]
    $ & $\displaystyle
    a \mathfrak{X}_1(z) + b \mathfrak{X}_2(z)
    $ \\\hline $
    x[n-1]
    $ & $\displaystyle
    z^{-1} \mathfrak{X}(z) + x[-1]
    $ \\\hline $\displaystyle
    x[n+1]
    $ & $\displaystyle
    z \mathfrak{X}(z) - z x[0]
    $ \\\hline $\displaystyle
    z_0^n x[n]
    $ & $\displaystyle
    \mathfrak{X}\qty(\frac{z}{z_0}) \vphantom{\mycases{1\\1}}
    $ \\\hline $\displaystyle
    x_{(k)}[n] = \mycases{
        x[n/k] \\ 0
    }
    $ & $\displaystyle
    \mathfrak{X}(z^k) \vphantom{\mycases{1\\1}}
    $ \\\hline $
    x^*[n]
    $ & $\displaystyle
    \mathfrak{X}^*(z^*)
    $ \\\hline $
    \begin{matrix}
        x_1[n] \ast x_2[n] \\
        \text{both causal}
    \end{matrix}
    $ & $\displaystyle
    \mathfrak{X}_1(z) \mathfrak{X}_2(z)
    $ \\\hline $\displaystyle
    n x[n]
    $ & $\displaystyle
    -z \dv{z} \mathfrak{X}(z) \vphantom{\mycases{1\\1}}
    $ \\\hline $\displaystyle
    \sum_{k=0}^{n} x[k]
    $ & $\displaystyle
    \frac{1}{1-z^{-1}} \mathfrak{X}(z) \vphantom{\mycases{1\\1}}
    $ \\\hline
\end{tabular}
\end{center}

\subsubsection{Initial- and Final-Value Theorem}
\begin{theorem}[Initial-Value and Final-Value Theorem]
    For any signal $x[n]$,
    \begin{gather*}
        x(0) = \lim_{z\to+\infty} \mathfrak{X}(z) \\
        \lim_{n\to+\infty} x[n] = \lim_{z\to 1} (z-1) \mathfrak{X}(z)
    \end{gather*}
\end{theorem}




\end{document}